\documentclass[11pt,a4paper]{article}
\usepackage[T1]{fontenc}
\usepackage{microtype}
\usepackage{isabelle,isabellesym}
\usepackage{amsmath}
\usepackage{amssymb}
\usepackage{booktabs}
\usepackage{tabularx}
\usepackage{url}

\newcolumntype{Y}{>{\raggedright\arraybackslash}X}
\newcommand{\idpath}[1]{\path{#1}}
\newcommand{\id}[1]{\ifmmode\mbox{\idpath{#1}}\else\idpath{#1}\fi}
\newcommand{\thy}[1]{\id{#1}}
\newcommand{\thm}[1]{\id{#1}}
\emergencystretch=3em

\isabellestyle{it}
\def\isacharunderscore{\discretionary{-}{}{-}}
\def\isacharunderscorekeyword{\discretionary{-}{}{-}}

\begin{document}

\title{Runtime Extensions for BMSSP in Isabelle/HOL}
\author{Arthur Freitas Ramos \and David Barros Hulak \and Ruy J. G. B. de Queiroz}
\maketitle

\begin{abstract}
This entry extends the AFP formalization
\emph{Correctness and Complexity of BMSSP in Isabelle/HOL}~\cite{BMSSP-AFP},
which formalizes the bounded multi-source shortest path algorithm of Duan,
Mao, Mao, Shu, and Yin~\cite{DuanMaoMaoShuYin2025}.  The published entry is
imported as a dependency and is not duplicated.

The extension adds a charged direct-insert relation that records both the
range charged by the recurrence and the semantic output returned by each
recursive child.  It proves that least costs are realized by verified-correct
runs, exhibits a closed branching instance, and records a
machine-checked totality obstruction for the older diagonal exact-concrete
path-family relation.  Its main theorem,
\texttt{path\_nonvac\_runtime\_bigo\_card\_V\_unconditional}, proves without
run-existence, cost-bound, or closed-bound assumptions that the least charged
cost on the growing unit-weight path family is
\(O(\lvert V_n\rvert \log^{2/3}\lvert V_n\rvert)\) when the internal schedule
is evaluated at \(n^2\).  This is a cost-relation theorem, not a machine-step
bound for the SML executable from the base entry.

AI assistance was used for proof engineering. The final definitions,
statements, and proofs are checked by Isabelle.
\end{abstract}

\tableofcontents

\section{Relationship to the Published Entry}

The base entry~\cite{BMSSP-AFP} supplies finite weighted digraph semantics,
BMSSP pre- and postconditions, the recursive and exact-concrete cost
relations, logarithmic schedule bounds, a bucketed partition model, and
executable correctness.  This entry is a child session of
\texttt{BMSSP\_Correctness}.  Its first three theories add only the new
charged relation and lemmas required by the follow-up; all other theories
import those extensions and the published AFP theories.

\begin{center}
\small
\begin{tabularx}{\textwidth}{@{}p{0.29\textwidth}Y@{}}
\toprule
\textbf{Theory} & \textbf{Added surface} \\
\midrule
\thy{BMSSP_Direct_Insert_Extensions} &
Charged direct-insert run and loop relations, correctness, existence, and
child-cost accounting. \\
\thy{BMSSP_Strict_Tie_Breaking_Extensions} &
Charged progress and level-bound closers under strict tie breaking. \\
\thy{BMSSP_Top_Level_Bounds_Extensions} &
Charged top-level run, cost, least-time, correctness, and refined-bound
interfaces. \\
\thy{BMSSP_Verified_Runtime} &
Joins the existing asymptotic least-cost theorem with correctness of a run
realizing that least cost. \\
\thy{BMSSP_Runtime_Instance_Branching} &
A closed out-star instance of the verified-runtime theorem. \\
\thy{BMSSP_Path_Family_Extended} &
The growing unit-weight path family, the unconditional cost-expression floor,
and conditional diagonal runtime statements. \\
\thy{BMSSP_Cost_Totality} &
The exact-concrete diagonal totality obstruction. \\
\thy{BMSSP_NonVacuous_Family} &
The decoupled charged family and unconditional non-vacuous runtime headline. \\
\bottomrule
\end{tabularx}
\end{center}

\section{The Unconditional Size-Parametric Headline}

For \(n\in\mathbb{N}\), let \(P_n\) be the directed unit-weight path with
vertices \(0,\ldots,n\).  The development proves
\begin{center}
\thm{path_nonvac_runtime_bigo_card_V_unconditional}
\end{center}
with the asymptotic meaning
\[
  T(n) \in
  O\!\left(\lvert V_n\rvert
    \bigl(\ln(\lvert V_n\rvert+2)\bigr)^{2/3}\right),
\]
where \(T(n)=\thm{path_nonvac_time}\ n\) is the least charged direct-insert
top-level cost for \(P_n\) at internal schedule \(n^2\).

The statement has no assumptions.  Its non-vacuity is explicit:
\thm{eventually_path_nonvac_charged_cost_exists} proves that charged runs
exist for all but finitely many \(n\), and
\thm{path_k_charged_cheap_imp_nonvac_time_le_bound} constructs a charged cost
below the refined graph-time budget.  The square schedule leaves the target
logarithmic class unchanged because \(\ln(n^2)=2\ln n\), while making the
level cap eventually large enough to cover the path.

\section{Verified Runtime and the Obstruction Record}

The theorem \thm{verified_bmssp_runtime_bigo_target} strengthens the base
entry's least-cost asymptotic statement by proving that the least cost is
eventually realized by an exact-concrete top-level run whose completed set is
the whole graph and whose output satisfies \thm{sssp_correct}.  The theory
\thy{BMSSP_Runtime_Instance_Branching} closes this statement on an out-star
whose source has two incomparable outgoing shortest paths.

Separately, \thy{BMSSP_Cost_Totality} proves the structural obstruction in the
older diagonal exact-concrete loop relation: a genuine step is forced to the
initial prefix under a sound label, and no genuine step can begin at a
positive lower bound.  The conditional diagonal family theorem is therefore
not used as evidence for non-vacuity.  The unconditional headline instead
uses the charged relation and decouples graph size from the internal schedule.

\section{Scope}

The imported base entry remains the source of executable
\texttt{bmssp\_distances} and its functional-correctness theorem.  This entry
does not claim that generated SML executes within the asymptotic cost bound.
The runtime statements concern the formal cost relations and their
graph-size-parametric instances.

\section{Checked Source}

\input{session}

\bibliographystyle{abbrv}
\bibliography{root}

\end{document}
